\begin{answers}{ch:training}{深いネットワークを学習させる技術}
\item 結果は $(-1.342, -0.447, 0.447, 1.342)$ で、$(2,4,6,8)$ のときとまったく同じです。
平均は $\mu=(1+3+5+7)/4=4$ で、平均との差は $(-3,-1,1,3)$ です。
これは $(2,4,6,8)$ の差と同じなので、分散も $\sigma^2=(9+1+1+9)/4=5$、$\sigma\approx2.236$ と同じになります。
$(1,3,5,7)$ は $(2,4,6,8)$ から全成分を 1 ずつ引いたものです。
全成分に同じ数を足し引きしても、平均が同じだけずれて差が変わらないので、出力は変わりません。
本文で見た「10 倍しても変わらない」と合わせると、レイヤー正規化は全体のずれにも倍率にも左右されず、成分どうしの相対的な並び方だけを残します。
\item 2 ステップとも $\hat m=2$、$\hat v=4$ で、2 ステップ目に動く量は $\eta$ の 1 倍（ちょうど $\eta$）です。
1 ステップ目は $m=0.1\times2=0.2$、$v=0.001\times4=0.004$ です。
補正すると $\hat m=0.2/0.1=2$、$\hat v=0.004/0.001=4$ になります。
2 ステップ目は $m=0.9\times0.2+0.1\times2=0.38$、$v=0.999\times0.004+0.001\times4=0.007996$ です。
補正の分母は $1-0.9^2=0.19$ と $1-0.999^2=0.001999$ なので、$\hat m=2$、$\hat v=4$ に戻ります。
動く量は $\eta\,\hat m/\sqrt{\hat v}=\eta\times2/2=\eta$ です（$\epsilon$ は無視しました）。
勾配が一定なら、補正後の $\hat m$ と $\hat v$ は毎回 $g$ と $g^2$ に一致します。
\item はい、最初の 1 歩は $(0.316, -0.316)$ で、学習率 0.1 の約 3.16 倍になりました（実測）。
$t=1$ では $m=0.1g$、$\sqrt{v}=\sqrt{0.001}\,\abs{g}$ なので、1 歩は $0.1/\sqrt{0.001}\approx3.16$ 倍の $\eta$ です。
勾配の大きさによらないので、$w_1$ も $w_2$ も同じ 0.316 だけ動きます。
その後も大きく動きすぎ、$w_2$ は 1 から 0.684、0.270 と進んで谷底を越え、$-0.68$ まで振れました。
100 歩後の損失は 0.0011 で、補正ありの 0.0007 より少し悪い値です。
\item 100 枚でも学習できます。\code{norm(h)} を \code{h} に変えると、100 枚では最初の損失が nan になります（実測）。
乱数の種 0・1・2 の 3 回で試すと、1,000 ステップ後の訓練損失は 3 回とも 0.0004 でした。
訓練の正解率は 100\%、テストは 99.3〜99.7\% です。
正規化を外すと、近道を通って足し込まれる量が 1 段ごとに倍々に増えます（\secref{sec:training-residual}）。
出力層の前で測った大きさは、種 0 で 10 枚で約 33、50 枚で約 $1.6\times10^{10}$、100 枚では約 $10^{20}$ でした。
50 枚までは最後の正規化がそろえ直すので、最初の損失は 1.1〜1.2 です。
100 枚では分散の計算で 2 乗した値が float32 の上限（約 $3.4\times10^{38}$）を超え、nan になります。
50 枚では発散こそしませんが、損失は $\log3=1.099$ から動かず、学習できませんでした（10 枚なら学習できます）。
\item 損失が最初に大きく跳ね上がり、最後まで上下に揺れます（実測。乱数の種 0・1・2 の平均）。
最初の 50 ステップの損失の最大は、元の設定（ウォームアップ 30、学習率 0.01）の 1.21 から 3.37 に増えました。
最後の 20 ステップの損失の平均も、0.0004 から 0.035 と約 100 倍に悪化しました。
訓練の正解率は 100\% から 99.0〜99.7\% に下がりましたが、損失は発散せず、学習は最後まで進みました。
片方だけ変えると、ウォームアップ 0・学習率 0.01 は最大 1.18、ウォームアップ 30・学習率 0.1 は 1.62 でした。
学習率が大きいときほど、最初の数ステップを小さな歩幅で始めるウォームアップが効くということです。
勾配クリッピングとレイヤー正規化があるので、この小さな問題では発散までは起きません。
\item 生成のときには、バッチの仲間がいないか、いても無関係な文なので、平均と分散がうまく求まらないからです。
学習時は多数の文をまとめて入れるので、ニューロンごとの平均と分散は多くの値から安定して求まります。
生成時は 1 本の文を 1 トークンずつ進めるので、その場で縦に平均をとると値は 1 個だけです。
平均は自分自身、分散は 0 になり、出力がすべて 0 になってしまいます。
そこで学習中にためた平均を使いますが、学習時と生成時で計算が変わり、不具合のもとになります。
また、複数の人の文を同時に生成すると、出力が無関係な仲間の文に左右されます。
レイヤー正規化は 1 トークンの $d$ 個の成分だけで決まるので、どちらの問題も起きません。
\end{answers}
